% Part: first-order-logic
% Chapter: syntax-and-semantics
% Section: first-order-languages

\documentclass[../../../include/open-logic-section]{subfiles}

\begin{document}

\olfileid{fol}{syn}{fol}

\olsection{لغات الرتبة الأولى}


أصل تعبيرات منطق الرتبة الأولى معجم يشتمل على
\emph{!!{variable}s} و\emph{!!{constant}s} و\emph{!!{predicate}s}،
وقد يشتمل أيضًا على \emph{!!{function}s}. فإذا ضُمّت إلى هذه الرموز
الروابط المنطقية والمكممات وعلامات الترقيم، من أقواس وفواصل، تألّفت
منها \emph{الحدود} و\emph{!!{formula}s}.

\begin{explain}
ووجه ذلك، قبل ضبطه بالتعريف، أن ال!!{predicate}s تسمّي الخصائص
والعلاقات، وال!!{constant}s تسمّي الكائنات المفردة، وال!!{function}s
تسمّي التطبيقات. وهذه الرموز، سوى ال!!{identity}~$\eq$، تسمى
\emph{الرموز غير المنطقية}، وباجتماعها تكون اللغة. فكل لغة من
الرتبة الأولى~$\Lang L$ إنما تتعيّن بتعيّن رموزها غير المنطقية.
وفي الحالة العامة يكون في $\Lang L$ من كل نوع رموز لا تتناهى.
\end{explain}

وتفصيل الرموز التي نعتمدها في الحالة العامة من منطق الرتبة الأولى كما يأتي:

\begin{enumerate}
\item الرموز المنطقية
\begin{enumerate}
\item الروابط المنطقية:
  \startycommalist
  \iftag{prvNot}{\ycomma $\lnot$ (النفي)}{}%
  \iftag{prvAnd}{\ycomma $\land$ (العطف)}{}%
  \iftag{prvOr}{\ycomma $\lor$ (الفصل)}{}%
  \iftag{prvIf}{\ycomma $\lif$ (!!{conditional})}{}%
  \iftag{prvIff}{\ycomma $\liff$ (!!{biconditional})}{}%
  \iftag{prvAll}{\ycomma $\lforall$ (المكمم الكلي)}{}%
  \iftag{prvEx}{\ycomma $\lexists$ (المكمم الوجودي)}{}.
\tagitem{prvFalse}{الثابت القضوي الدال على !!{falsity}~$\lfalse$.}{}
\tagitem{prvTrue}{الثابت القضوي الدال على !!{truth}~$\ltrue$.}{}
\item ال!!{identity} الثنائية المحل~$\eq$.
\item مجموعة !!{denumerable} من ال!!{variable}s: $\Obj v_0$، و$\Obj v_1$، و$\Obj
  v_2$، و\dots
\end{enumerate}
\item الرموز غير المنطقية التي تؤلف \emph{اللغة القياسية} لمنطق
  الرتبة الأولى
\begin{enumerate}
\item مجموعة !!{denumerable} من ال!!{predicate}s ذات $n$ من المحال لكل
  $n>0$: $\Obj A^n_0$، و$\Obj A^n_1$، و$\Obj A^n_2$، و\dots
\item مجموعة !!{denumerable} من ال!!{constant}s: $\Obj c_0$، و$\Obj c_1$، و$\Obj
  c_2$، و\dots.
\item مجموعة !!{denumerable} من ال!!{function}s ذات $n$ من المحال لكل
  $n>0$: $\Obj f^n_0$، و$\Obj f^n_1$، و$\Obj f^n_2$، و\dots
\end{enumerate}
\item علامات الترقيم: القوسان ( و)، والفاصلة.
\end{enumerate}

وعلى اللغة القياسية الكاملة لمنطق الرتبة الأولى تُحمل أكثر
تعريفاتنا ونتائجنا. وقد يقتضي موضع الاستعمال أن نقتصر من
ال!!{predicate}s وال!!{constant}s وال!!{function}s على قليل.

\begin{ex}
فمن ذلك لغة الحساب~$\Lang L_A$؛ إذ ليس فيها إلا !!{predicate}
ثنائي المحل~$<$، و!!{constant} واحد~$\Obj 0$، و!!{function}
أحادية المحل~$\prime$، واثنتان من ال!!{function}s الثنائية المحل، هما~$+$ و~$\times$.
\end{ex}

\begin{ex}
وأما لغة نظرية المجموعات~$\Lang L_Z$، فليس من رموزها غير المنطقية
إلا ال!!{predicate} الثنائي المحل~$\in$.
\end{ex}

\begin{ex}
وكذلك لغة الترتيبات~$\Lang L_\le$، فإن رموزها غير المنطقية تنحصر في
ال!!{predicate} الثنائي المحل~$\le$.
\end{ex}

وليس في اختيار هذه الرسوم إلا اصطلاح في التسمية؛ فهي، عند الضبط
الصوري، أسماء مستعارة. فمثلًا، $<$
و$\in$ و$\le$ أسماء مستعارة لـ~$\Obj A^2_0$، و$\Obj 0$ اسم مستعار
لـ~$\Obj c_0$، و$\prime$ لـ~$\Obj f^1_0$، و$+$ لـ~$\Obj f^2_0$،
و$\times$ لـ~$\Obj f^2_1$.

\iftag{defNot,defOr,defAnd,defIf,defIff,defTrue,defFalse,defEx,defAll}{%
ومع الروابط الأولية و
\iftag{notprvEx,notprvAll}{المكمم}{المكممات} التي تقدمت، نأخذ
الرموز الآتية على أنها رموز \emph{معرّفة}:
\startycommalist
  \iftag{defNot}{\ycomma $\lnot$ (النفي)}{}%
  \iftag{defAnd}{\ycomma $\land$ (العطف)}{}%
  \iftag{defOr}{\ycomma $\lor$ (الفصل)}{}%
  \iftag{defIf}{\ycomma $\lif$ (!!{conditional})}{}%
  \iftag{defIff}{\ycomma $\liff$ (!!{biconditional})}{}%
  \iftag{defAll}{\ycomma $\lforall$ (المكمم الكلي)}{}%
  \iftag{defEx}{\ycomma $\lexists$ (المكمم الوجودي)}{}%
  \iftag{defFalse}{\ycomma !!{falsity}~$\lfalse$}{}%
  \iftag{defTrue}{\ycomma !!{truth}~$\ltrue$}{}%
}{}.

\begin{tagblock}{defNot,defOr,defAnd,defIf,defIff,defTrue,defFalse,defEx,defAll}
\begin{explain}
أما الرمز المعرّف فليس من أصل اللغة في اصطلاحها الصوري، وإنما هو
اختصار نضعه للتيسير. فلولاه، ولزمنا الاقتصار على الرموز الأولية،
لطالت بعض الصيغ طولًا شديدًا. وفوق اختصار الصيغ فائدة أخرى، هي
اختصار البراهين؛ وذلك أن لكل رمز أولي حالةً تخصه في البرهان،
فكلما ازداد عدد تلك الرموز ازدادت الحالات وطال البرهان.
\end{explain}
\end{tagblock}

% Alternate symbols

\begin{intro}
وقد تختلف اصطلاحات الكتب والمدرّسين عما اخترناه هنا. فمن الرموز
الشائعة عندهم $\sim$ و$\neg$ و~!\ للدلالة على
«النفي»، واستعمال $\wedge$ و$\cdot$ و$\&$ للدلالة على «العطف». أما
الرموز الشائعة لـ«الشرط» أو «الاستلزام» فهي $\rightarrow$
و$\Rightarrow$ و$\supset$.
\iftag{prvIff,defIff}{أما رموز «الشرط الثنائي» أو «الاستلزام المتبادل»
  أو «التكافؤ (المادي)» فهي $\leftrightarrow$ و$\Leftrightarrow$
  و$\equiv$.}{}
\iftag{prvFalse,defFalse}{تتنوع أسماء الرمز~$\lfalse$؛ فيسمى
  «الكذب» أو «الباطل» أو «المحال» أو «القاع».}{}
\iftag{prvTrue,defTrue}{تتنوع أسماء الرمز~$\ltrue$؛ فيسمى
  «الصدق» أو «الحق» أو «القمة».}{}

ومن اصطلاحهم أيضًا أن يجعلوا الأحرف الصغيرة من صدر الأبجدية اللاتينية (مثل
$a$ و$b$ و$c$) لل!!{constant}s (وتسمى أحيانًا أسماء)، والأحرف
الصغيرة من آخرها (مثل $x$ و$y$ و$z$) لل!!{variable}s. وتقترن
المكممات بال!!{variable}s، مثل~$x$؛ ومن الاختلافات الترميزية
$\forall x$ و$(\forall x)$ و$(x)$ و$\Pi x$ و$\bigwedge_x$ للمكمم
الكلي، و$\exists x$ و$(\exists x)$ و$(Ex)$ و$\Sigma x$ و
$\bigvee_x$ للمكمم الوجودي.
\end{intro}

\begin{explain}
ولنا أن نجعل المؤثرات القضوية كلها، ومعها المكممان، رموزًا أولية؛
ولنا أن نكتفي بأصل أقل عددًا، ثم نعرّف سائر
ال!!{operator}s به. فمن مجموعات المؤثرات البوليانية التي تكون
«كاملة من حيث الدوال الصدقية» $\{ \lnot, \lor \}$ و
$\{ \lnot, \land \}$ و$\{ \lnot, \lif\}$؛ ويمكن ضم كل منها إلى أي
من المكممين للحصول على لغة رتبة أولى كاملة القدرة التعبيرية.

ومن ال!!{operator}s الأخرى شرطة شيفر~$|$، ونسبتها إلى
هنري شيفر، وسهم بيرس~$\downarrow$، ويقال له أيضًا خنجر كواين.
فإذا أريد بهما معناهما المعتاد، أي «نفي العطف» للأول و«نفي الفصل»
للثاني، استقل كل واحد منهما بالكمال من حيث الدوال الصدقية.
\end{explain}

\end{document}
